\scanpage{10}

% loc: top, abandoned opening
Soit $A$ une \deleted{$\unsure{\otimes}$}-catégorie \deleted{\unsure{commutative}}\\
\deleted{\unsure{Comme une|une} $\otimes$ catégorie $\unsure{\mathrm{AUC}|\mathrm{ANC}}$ \unsure{est triviale|dont tous} (et définie à isomorphisme \unsure{)}}

% loc: upper paragraph
\deleted{\uline{\unsure{Problème|Définition}}} Soit $\Phi(A)$ la catégorie des familles finies d'objets de $A$, un morphisme de $(I,(L_i)_{i\in I})$ dans $(J,(M_j)_{j\in J})$ étant un couple $(\tau,(\varphi_i)_{i\in I})$ d'une bijection $\tau:I\to J$ et d'isom. $\varphi_i:L_i\xrightarrow{\sim}M_{\tau(i)}$. On introduit sur $\Phi(A)$ une $\otimes$-structure en posant

% loc: tensor formula
\[(I,(L_i)_{i\in I})\otimes(J,(M_j)_{j\in J})=(I\amalg J,(L_i;M_j))\ \ldots\]

% loc: structural constraints
\deleted{\unsure{et}} une structure \deleted{d'associativité}\inserted{d'associativité}, provenant de l'associativité dans l'opération $I\amalg J$, une structure unité par la famille vide \unsure{étant|et avec|étant aussi} l'objet unité, une structure de commutativité par la commutativité de $I\amalg J$. (NB. pour ces définitions, la $\otimes$-structure de $C$ n'a pas encore servi.)

% loc: cancelled functor draft
On \deleted{peut trouver un $\otimes$-foncteur AUC} \deleted{\unsure{pour}}
\[
\deleted{\Pi:\Phi(C)\longrightarrow C\qquad
\substack{\Phi_n(C)=\text{ss-catégorie}\\\text{pleine de }\Phi(C)\text{ formée}}}
\]
\deleted{tel que (la restriction de $\Pi$ au \unsure{le}}\inserted{\deleted{\unsure{notant} $\Pi_n$}}

% loc: canonical functor
a un foncteur canonique
\[\Phi(A)\longrightarrow\big(\underset{\substack{\text{avec isom}\\\text{comme}\\\text{morphismes}}}{\mathrm{Ens\ finis}}\big)\deleted{\unsure{\text{ds Ens}}}\]

% loc: middle, cardinality
(foncteur transportable \ldots), on désigne par $\Phi_n(A)$ l'image inverse \deleted{\unsure{de la}} la sous-catégorie des ensembles de \deleted{\unsure{dim}} cardinal $n$ ($n$ entier $\geq 0$). On a donc

% loc: middle display
\[\Phi(A)=\coprod_n\Phi_n(A),\quad\left\{\begin{aligned}
\Phi_0(A)&=\{e\}\\
\Phi_1(A)&\simeq A
\end{aligned}\right.\]

% loc: 1, statement
1) Soit $C$ une catégorie munie d'un produit tensoriel $\otimes$ AUC. Considérons le foncteur ``restriction à $\Phi_1(A)$''

% loc: restriction display
\[\underline{\operatorname{Hom}}^{\otimes\,\mathrm{AUC}}(\Phi(A),C)\longrightarrow\underline{\operatorname{Hom}}(A,C)\]

% loc: proposition
\uline{Proposition}\quad C'est une équivalence de catégories.

% loc: proof
\uline{Dém} \ldots.

% loc: corollary
\uline{Corollaire}\quad \inserted{Soit $C$ une $\otimes$-catégorie AUC.} Il y a, à isom unique près, un et un seul $\otimes$-foncteur AUC

% loc: bottom
\[\Pi:\Phi(e)\longrightarrow C\qquad\Pi=(\Pi_n)\]

% archive sheet number: 9
